% =========================================================
% BÁO CÁO KỸ THUẬT — BÀI TẬP LỚN XỬ LÝ TÍN HIỆU SỐ
% Chủ đề 4: Hệ thống tự động phát và giải mã tín hiệu DTMF
%
% Biên dịch:  pdflatex main  ->  biber main  ->  pdflatex main x2
% Hoặc:       latexmk -pdf main.tex
%
% QUY TẮC: mỗi người chỉ sửa file chương của mình trong thư mục Chapters/.
% Không ai sửa dtmf_report.cls và phần preamble dưới đây.
% =========================================================
\documentclass{dtmf_report}

\addbibresource{references.bib}

% ---- Thông tin trang bìa: sửa cho đúng trường, lớp, giảng viên ----
\settruong{HỌC VIỆN AN NINH NHÂN DÂN}
\setkhoa{KHOA AN NINH MẠNG VÀ PHÒNG, CHỐNG TỘI PHẠM SỬ DỤNG CÔNG NGHỆ CAO}
\sethocphan{XỬ LÝ TÍN HIỆU SỐ}
\setdetai{Xây dựng hệ thống tự động phát và giải mã\\ tín hiệu điện thoại DTMF}
\setnhom{Nhóm 4}
\setgiangvien{TS. Nguyễn Văn A}
\setdiadiem{Hà Nội}
\setthoigian{9/2026}
\setthanhvien{
  \renewcommand{\arraystretch}{1.12}%
  \begin{tabular}{@{}r@{~}l l @{\hspace{13pt}} r@{~}l l@{}}
    1. & Đỗ Hồng Nhất      & B9D55  & 10. & Nguyễn Tiến Mạnh        & B9D55 \\
    2. & Nguyễn Thu Hà     & B8D55  & 11. & Nguyễn Huy Hoàng        & B9D55 \\
    3. & Lê Tuấn Anh       & B9D55  & 12. & Nguyễn Cửu An Phú       & B9D55 \\
    4. & Nguyễn Quốc Đăng  & B9D55  & 13. & Đinh Minh Trí           & B9D55 \\
    5. & Lê Đức Anh        & B9D55  & 14. & Đặng Nguyễn Lâm         & B9D55 \\
    6. & Nguyễn Minh Hiếu  & B9D55  & 15. & Huỳnh Nguyễn Bảo Khương & B9D55 \\
    7. & Hồ Văn Hoàng Quân & B10D55 & 16. & Nguyễn Duy Tiến         & B9D55 \\
    8. & Nguyễn Trí Dũng   & B9D55  & 17. & Cao Tiến Thành          & B9D55 \\
    9. & Trương Minh Tân   & B9D55  & 18. & Nguyễn Sỹ Long          & B9D55 \\
  \end{tabular}
}

\begin{document}

% =========================================================
% PHẦN ĐẦU
% =========================================================
\trangbia

\pagenumbering{roman}

\tableofcontents
\clearpage

\listoffigures
\clearpage

\listoftables
\clearpage

\unsection{DANH MỤC TỪ VIẾT TẮT VÀ KÝ HIỆU}

\begin{table}[H]
\centering
\small
\begin{tabular}{l l}
\toprule
\textbf{Ký hiệu} & \textbf{Ý nghĩa} \\
\midrule
DTMF   & Dual-Tone Multi-Frequency — âm đa tần kép \\
DFT    & Discrete Fourier Transform — biến đổi Fourier rời rạc \\
FFT    & Fast Fourier Transform — biến đổi Fourier nhanh \\
FIR    & Finite Impulse Response — đáp ứng xung hữu hạn \\
IIR    & Infinite Impulse Response — đáp ứng xung vô hạn \\
BIBO   & Bounded Input Bounded Output — vào giới hạn, ra giới hạn \\
SNR    & Signal-to-Noise Ratio — tỷ số tín hiệu trên nhiễu \\
ITU-T  & International Telecommunication Union — Telecom Standardization \\
\midrule
$\fs$        & Tần số lấy mẫu \\
$N$          & Độ dài khung phân tích (số mẫu) \\
$k$          & Chỉ số bin tần số trong DFT \\
$\dfres$     & Độ phân giải tần số, $\dfres = \fs/N$ \\
$\wzero$     & Tần số cộng hưởng chuẩn hoá \\
$r$          & Bán kính điểm cực của bộ lọc cộng hưởng \\
$\Qfac$      & Hệ số phẩm chất của bộ lọc \\
$\Htf$       & Hàm truyền của hệ thống trong miền $z$ \\
\bottomrule
\end{tabular}
\end{table}
\clearpage

% =========================================================
% NỘI DUNG CHÍNH
% =========================================================
\pagenumbering{arabic}

% ---------------------------------------------------------
% MỞ ĐẦU  (không đánh số chương)
% ---------------------------------------------------------
\unsection{MỞ ĐẦU}

\unsubsection{1. Đặt vấn đề}
% Vai trò của báo hiệu đa tần trong mạng điện thoại. Bài toán: từ một đoạn
% âm thanh chứa các phím bấm, tự động khôi phục chuỗi ký tự đã bấm.
% Ứng dụng: tổng đài tự động, hệ thống IVR, điều khiển thiết bị qua đường thoại.
[Nội dung mục 1 viết tại đây.]

\unsubsection{2. Mục tiêu đề tài}
% Bốn mục tiêu, phát biểu ngắn gọn.
\begin{enumerate}[leftmargin=*]
  \item Bộ phát tín hiệu DTMF đúng chuẩn ITU-T, có mô phỏng nhiễu.
  \item Ba bộ giải mã độc lập: DFT/FFT, Goertzel, ngân hàng bộ lọc.
  \item So sánh định lượng ba phương pháp về độ chính xác, độ bền nhiễu và chi phí tính toán.
  \item Phần mềm mô phỏng có giao diện trực quan.
\end{enumerate}

\unsubsection{3. Phạm vi, giới hạn và phương pháp thực hiện}
% 12 phím chuẩn, không xử lý cột 1633 Hz; fs = 8 kHz cố định; xử lý theo khối;
% dữ liệu tổng hợp có nhãn làm chính, ghi âm thực tế để kiểm chứng bổ sung.
[Nội dung mục 3 viết tại đây.]

\clearpage

% ---------------------------------------------------------
% I. CƠ SỞ LÝ THUYẾT
% ---------------------------------------------------------
\section{CƠ SỞ LÝ THUYẾT}

\subsection{Tín hiệu DTMF}

\subsubsection{Nguyên lý âm kép và bảng tần số chuẩn}
% Mỗi phím phát đồng thời một tần số nhóm thấp (hàng) và một tần số nhóm cao (cột);
% ma trận 4 x 3 tạo 12 tổ hợp.
[Nội dung viết tại đây.]

\begin{table}[H]
\centering
\caption{Bảng tần số chuẩn DTMF theo khuyến nghị ITU-T Q.23}
\label{tab:tanso}
\begin{tabular}{l c c c}
\toprule
& \textbf{1209 Hz} & \textbf{1336 Hz} & \textbf{1477 Hz} \\
\midrule
\textbf{697 Hz} & 1 & 2 & 3 \\
\textbf{770 Hz} & 4 & 5 & 6 \\
\textbf{852 Hz} & 7 & 8 & 9 \\
\textbf{941 Hz} & * & 0 & \# \\
\bottomrule
\end{tabular}
\end{table}

\subsubsection{Cơ sở lựa chọn bộ bảy tần số}
% MỤC PHÂN TÍCH - VIẾT KỸ. Ba tiêu chí:
%  (a) không tần số nào là bội số nguyên của tần số khác;
%  (b) hai nhóm tách biệt đủ xa để lọc riêng;
%  (c) toàn dải nằm trong băng thoại 300-3400 Hz.
[Nội dung viết tại đây.]

\subsubsection{Mô hình toán học và tham số chuẩn}
% Biểu thức tín hiệu liên tục và rời rạc; tham số ITU-T Q.23/Q.24.
Tín hiệu DTMF ứng với một phím là tổng của hai dao động hình sin:
\begin{equation}\label{eq:dtmf-lien-tuc}
  x(t) = A_{1}\sin(2\pi f_{1} t) + A_{2}\sin(2\pi f_{2} t)
\end{equation}
Sau khi lấy mẫu ở tần số $\fs$, tín hiệu rời rạc có dạng
\begin{equation}\label{eq:dtmf-roi-rac}
  x[n] = A_{1}\sin\!\left(\frac{2\pi f_{1} n}{\fs}\right)
       + A_{2}\sin\!\left(\frac{2\pi f_{2} n}{\fs}\right)
\end{equation}

\subsection{Lấy mẫu và hiện tượng chồng phổ}
% Định lý Nyquist-Shannon; kiểm chứng cho fmax = 1477 Hz với fs = 8 kHz;
% cơ chế chồng phổ kèm mô phỏng minh hoạ.
[Nội dung viết tại đây.]

\chohinh{minh hoạ hiện tượng chồng phổ}

\subsection{Phổ tần số của tín hiệu DTMF}
% Phổ biên độ của một phím - hai đỉnh rõ rệt. Đặc điểm then chốt: DTMF lý tưởng
% gần như không có hài bậc cao, khác hẳn tiếng nói -> cơ sở cho kiểm tra
% chống talk-off ở mục 2.2.4. Ảnh hưởng của ba loại nhiễu lên phổ.
[Nội dung viết tại đây.]

\clearpage

% ---------------------------------------------------------
% II. CÁC PHƯƠNG PHÁP GIẢI MÃ
% ---------------------------------------------------------
\section{CÁC PHƯƠNG PHÁP GIẢI MÃ}

\subsection{Phương pháp DFT/FFT}

\subsubsection{Biến đổi Fourier rời rạc}
% Công thức DFT, ý nghĩa X[k], quan hệ f_k = k*fs/N, phổ công suất,
% tính đối xứng với tín hiệu thực; độ phức tạp O(N^2) so với O(N log N).
Với một khung $N$ mẫu, hệ số DFT tại bin thứ $k$ được tính bằng
\begin{equation}\label{eq:dft}
  X[k] = \sum_{n=0}^{N-1} x[n]\, e^{-j 2\pi k n / N},
  \qquad k = 0, 1, \dots, N-1
\end{equation}
Bin thứ $k$ tương ứng với tần số thực $f_{k} = k\,\fs / N$.

\subsubsection{Độ phân giải tần số và đánh đổi thời gian--tần số}
% MỤC BIỆN MINH THIẾT KẾ - VIẾT KỸ.
% Delta f = fs/N; khoảng cách nhỏ nhất cần phân biệt là 73 Hz (697 và 770 Hz);
% vì sao không chọn N lớn hơn -> kết luận chọn khung 256 điểm.
Độ phân giải tần số của phép biến đổi là
\begin{equation}\label{eq:do-phan-giai}
  \dfres = \frac{\fs}{N}
\end{equation}
[Nội dung viết tại đây.]

\subsubsection{Rò rỉ phổ và hàm cửa sổ}
% Nguyên nhân rò rỉ; búp chính và búp phụ; bảng so sánh bốn hàm cửa sổ;
% lập luận chọn Hamming.
[Nội dung viết tại đây.]

\begin{table}[H]
\centering
\caption{So sánh các hàm cửa sổ}
\label{tab:cua-so}
\begin{tabular}{l c c l}
\toprule
\textbf{Hàm cửa sổ} & \textbf{Độ rộng búp chính} & \textbf{Mức búp phụ} & \textbf{Nhận xét} \\
\midrule
Chữ nhật &  &  &  \\
Hamming  &  &  &  \\
Hanning  &  &  &  \\
Blackman &  &  &  \\
\bottomrule
\end{tabular}
\end{table}

\subsubsection{Áp dụng vào giải mã DTMF}
% Phân khung 256 điểm, bước nhảy 128, chồng lấp 50%; ánh xạ 7 tần số chuẩn
% sang chỉ số bin; thuật toán dò đỉnh trong hai nhóm.
[Nội dung viết tại đây.]

\chohinh{phổ FFT của phím "5" có đánh dấu 8 bin chuẩn}

\subsection{Phương pháp thuật toán Goertzel}

\subsubsection{Đặt vấn đề}
% Bài toán chỉ cần năng lượng tại 8 tần số xác định trước;
% chi phí lãng phí khi tính đủ 256 bin nhưng chỉ dùng 8.
[Nội dung viết tại đây.]

\subsubsection{Dẫn giải công thức từ DFT}
% SÁU BƯỚC, KHÔNG NHẢY BƯỚC. Đây là phần toán trọng tâm của cả báo cáo.
\begin{enumerate}[label=\textbf{Bước \arabic*.}, leftmargin=*,
                  labelsep=0.5em, itemsep=8pt, topsep=8pt, parsep=4pt]

  \item Viết DFT dưới dạng tích luỹ đệ quy. Đặt $W = e^{j2\pi k/N}$ và định nghĩa
        chuỗi trung gian
        \begin{equation}\label{eq:goertzel-b1}
          y[n] = x[n] + W\, y[n-1], \qquad y[-1] = 0
        \end{equation}

  \item {}[Chứng minh $y[N-1]$ cho ra đúng $\Xk$ sau $N$ bước lặp.]

  \item {}[Nhận xét bộ lọc bậc một này có hệ số phức, chi phí tính toán cao.]

  \item {}[Ghép hàm truyền với liên hợp phức của chính nó để khử phần ảo.]

  \item Phương trình sai phân bậc hai hệ số thực thu được:
        \begin{equation}\label{eq:goertzel-b5}
          s[n] = x[n] + 2\cos\!\left(\frac{2\pi k}{N}\right) s[n-1] - s[n-2]
        \end{equation}

  \item Công suất tại bin $k$ sau $N$ mẫu:
        \begin{equation}\label{eq:goertzel-b6}
          P = s[N-1]^{2} + s[N-2]^{2}
              - 2\cos\!\left(\frac{2\pi k}{N}\right) s[N-1]\, s[N-2]
        \end{equation}
\end{enumerate}

\subsubsection{Goertzel dưới góc nhìn bộ lọc số}
% Hàm truyền H(z) của bộ cộng hưởng bậc hai; sơ đồ khối dạng trực tiếp II;
% vị trí cặp cực trên đường tròn đơn vị; vì sao không mất ổn định.
% Nhận xét cốt lõi: mỗi mẫu chỉ cần một phép nhân thực.
[Nội dung viết tại đây.]

\chohinh{sơ đồ khối bộ lọc Goertzel bậc hai}

\subsubsection{Lựa chọn tham số và luật quyết định}
% Tiêu chí chọn N; bảng bin k với N = 205 kèm % sai lệch;
% luật quyết định bốn tầng; cơ chế gộp khung liên tiếp.
[Nội dung viết tại đây.]

\begin{table}[H]
\centering
\caption{Chỉ số bin $k$ và sai lệch tần số với $N = 205$, $\fs = 8000$ Hz}
\label{tab:bin-k}
\begin{tabular}{c c c c}
\toprule
\textbf{$f$ chuẩn (Hz)} & \textbf{$k$} & \textbf{$f$ của bin (Hz)} & \textbf{Sai lệch} \\
\midrule
697  & 18 & 702{,}44  & $+0{,}78\%$ \\
770  & 20 & 780{,}49  & $+1{,}36\%$ \\
852  & 22 & 858{,}54  & $+0{,}77\%$ \\
941  & 24 & 936{,}59  & $-0{,}47\%$ \\
1209 & 31 & 1209{,}76 & $+0{,}06\%$ \\
1336 & 34 & 1326{,}83 & $-0{,}69\%$ \\
1477 & 38 & 1482{,}93 & $+0{,}40\%$ \\
\bottomrule
\end{tabular}
\end{table}

\begin{notebox}[Luật quyết định bốn tầng]
\begin{enumerate}[leftmargin=*,nosep]
  \item \textbf{Ngưỡng đỉnh} — đỉnh mạnh nhất mỗi nhóm vượt đỉnh thứ nhì cùng nhóm ít nhất $6\dB$.
  \item \textbf{Tỷ lệ năng lượng} — tổng năng lượng 8 bin chuẩn chiếm $\geq 70\%$ năng lượng khung.
  \item \textbf{Kiểm tra twist} — chênh lệch biên độ hai nhóm nằm trong giới hạn cho phép.
  \item \textbf{Kiểm tra hài bậc hai} — năng lượng tại $2f_{\text{hàng}}$ và $2f_{\text{cột}}$ phải thấp.
\end{enumerate}
\end{notebox}

\subsubsection{Phân tích chi phí tính toán}
% Goertzel 8x205 ~ 1640; FFT 256 radix-2 ~ 4096; ngân hàng 8 IIR ~ 3280.
% Ghi rõ TỪNG CON SỐ RA TỪ ĐÂU. Phân tích ngưỡng hoà vốn.
[Nội dung viết tại đây.]

\subsubsection{Kiểm chứng thuật toán}
% Ví dụ N = 16, k = 3, x[n] = cos(2*pi*3n/16) -> P = 64,000000.
% Đối chiếu lý thuyết X[3] = N/2 = 8 nên |X[3]|^2 = 64.
% Đối chiếu thêm với hàm goertzel() của MATLAB.
[Nội dung viết tại đây.]

\begin{matlabbox}[Kiểm chứng \texttt{goertzel\_power.m}]
n = 0:15;
x = cos(2*pi*3*n/16);
P = goertzel_power(x, 3, 16);   % ky vong: P = 64.000000
\end{matlabbox}

\subsection{Phương pháp ngân hàng bộ lọc thông dải}

\subsubsection{Nguyên lý và lựa chọn họ bộ lọc}
% Tín hiệu qua song song 8 bộ lọc cộng hưởng, tách đường bao, so ngưỡng.
% Bảng so sánh FIR/IIR; lập luận chọn IIR.
[Nội dung viết tại đây.]

\begin{table}[H]
\centering
\caption{So sánh bộ lọc FIR và IIR}
\label{tab:fir-iir}
\begin{tabular}{l p{5cm} p{5cm}}
\toprule
& \textbf{FIR} & \textbf{IIR} \\
\midrule
Hồi tiếp     & Không & Có \\
Tính ổn định & Luôn ổn định & Cần kiểm tra vị trí điểm cực \\
Số hệ số     & Nhiều hơn & Ít hơn \\
Đặc tính pha & Có thể tuyến tính tuyệt đối & Thường không tuyến tính \\
\bottomrule
\end{tabular}
\end{table}

\subsubsection{Thiết kế bằng biến đổi song tuyến}
% Nguyên mẫu Butterworth; công thức biến đổi; chứng minh ánh xạ bảo toàn ổn định;
% méo tần số và tiền méo.
Phép biến đổi song tuyến ánh xạ mặt phẳng $s$ sang mặt phẳng $z$ theo
\begin{equation}\label{eq:song-tuyen}
  s = \frac{2}{T}\cdot\frac{1 - z^{-1}}{1 + z^{-1}}
\end{equation}
kéo theo quan hệ méo tần số
\begin{equation}\label{eq:pre-warp}
  \Omega = \frac{2}{T}\tan\!\left(\frac{\omega}{2}\right)
\end{equation}
Phương án cộng hưởng trực tiếp cho mỗi tần số chuẩn:
\begin{equation}\label{eq:cong-huong}
  \Htf = G\cdot\frac{1 - z^{-2}}
                    {1 - 2r\cos(\wzero) z^{-1} + r^{2} z^{-2}}
\end{equation}

\subsubsection{Điểm cực, điểm không và ổn định BIBO}
% Điều kiện |pole| < 1; giản đồ cực-không của 8 bộ lọc;
% chứng minh hệ ổn định với r = 0,99.
[Nội dung viết tại đây.]

\chohinh{giản đồ cực--không của 8 bộ lọc thông dải}

\subsubsection{Băng thông và hệ số phẩm chất Q}
% BW ~ (1-r)*fs/pi; Q = f0/BW; bảng Q cho 8 bộ lọc;
% vì sao Q tăng theo tần số; đánh đổi ba chiều.
Với $r$ tiến gần 1, băng thông và hệ số phẩm chất xấp xỉ bằng
\begin{equation}\label{eq:bw-q}
  \mathrm{BW} \approx \frac{(1-r)\,\fs}{\pi},
  \qquad
  \Qfac = \frac{f_{0}}{\mathrm{BW}}
\end{equation}

\subsubsection{Ảnh hưởng của lượng tử hoá hệ số}
% Dấu phẩy động 64-bit so với dấu phẩy cố định 16-bit;
% cơ chế dịch chuyển điểm cực; vì sao bộ lọc Q cao nhạy cảm hơn;
% nguy cơ cực bị đẩy ra ngoài đường tròn đơn vị.
[Nội dung viết tại đây.]

\subsection{So sánh ba phương pháp về mặt lý thuyết}
% Bảng đối chiếu cơ sở toán học, độ phân giải, chi phí, độ trễ,
% độ phức tạp cài đặt. Nhận xét sơ bộ trước khi kiểm chứng ở Chương IV.
[Nội dung viết tại đây.]

\clearpage

% ---------------------------------------------------------
% III. THIẾT KẾ VÀ XÂY DỰNG PHẦN MỀM
% ---------------------------------------------------------
\section{THIẾT KẾ VÀ XÂY DỰNG PHẦN MỀM}

\subsection{Kiến trúc hệ thống}
% Sơ đồ phân lớp: Giao diện -> Điều phối -> ba bộ giải mã -> Tiện ích -> Dữ liệu.
% Nguyên tắc tách bạch trách nhiệm. Lưu đồ một lượt xử lý đầy đủ.
[Nội dung viết tại đây.]

\chohinh{sơ đồ kiến trúc phần mềm theo mô-đun}

\chohinh{lưu đồ luồng xử lý một lần bấm phím}

\subsection{Tổ chức mã nguồn và chuẩn giao tiếp}
% Cấu trúc thư mục; bảng danh sách hàm chính;
% ba bộ giải mã dùng chung một chữ ký hàm.
[Nội dung viết tại đây.]

\dirtree{%
.1 DMTF/.
.2 src/.
.3 gen/ \qquad \textit{sinh tín hiệu và nhiễu}.
.3 decode/ \qquad \textit{ba bộ giải mã}.
.3 util/ \qquad \textit{phân khung, quyết định, đánh giá}.
.2 app/ \qquad \textit{giao diện App Designer}.
.2 tests/ \qquad \textit{kiểm thử tự động}.
.2 data/ \qquad \textit{dữ liệu và hệ số bộ lọc}.
.2 results/ \qquad \textit{hình và số liệu benchmark}.
}

\begin{table}[H]
\centering
\caption{Danh sách các hàm chính}
\label{tab:ham-chinh}
\small
\begin{tabular}{@{}l p{2.7cm} p{2.7cm} p{3.4cm}@{}}
\toprule
\textbf{Hàm} & \textbf{Đầu vào} & \textbf{Đầu ra} & \textbf{Chức năng} \\
\midrule
\texttt{dtmf\_generate}   & & & \\
\texttt{dtmf\_addnoise}   & & & \\
\texttt{dtmf\_segment}    & & & \\
\texttt{dtmf\_decode\_fft}& & & \\
\texttt{goertzel\_power}  & & & \\
\texttt{design\_bpf\_bank}& & & \\
\texttt{dtmf\_decide}     & & & \\
\texttt{dtmf\_metrics}    & & & \\
\bottomrule
\end{tabular}
\end{table}

\subsection{Giao diện người dùng}
% Ba tab: Giải mã, Bộ lọc, Đánh giá. Bảng điều khiển.
% Lý giải thiết kế biểu đồ 8 thanh năng lượng.
[Nội dung viết tại đây.]

\chohinh{ảnh chụp ba tab giao diện}

\clearpage

% ---------------------------------------------------------
% IV. KẾT QUẢ THỰC NGHIỆM VÀ ĐÁNH GIÁ
% ---------------------------------------------------------
\section{KẾT QUẢ THỰC NGHIỆM VÀ ĐÁNH GIÁ}

\subsection{Thiết lập thí nghiệm và tiêu chí đánh giá}
% Môi trường; bộ dữ liệu kiểm thử; mức SNR từ -10 đến 30 dB; ba loại nhiễu.
% Tiêu chí: độ chính xác, khoảng cách Levenshtein, ma trận nhầm lẫn, SNR, thời gian.
Tỷ số tín hiệu trên nhiễu được tính theo
\begin{equation}\label{eq:snr}
  \SNRdB = 10\log_{10}\!
  \left(\frac{\sum_{n} x^{2}[n]}{\sum_{n}\big(y[n]-x[n]\big)^{2}}\right)
\end{equation}

\subsection{Kết quả trên dữ liệu sạch}
% Độ chính xác cả ba phương pháp; phân tích sai sót còn lại.
[Nội dung viết tại đây.]

\subsection{Độ bền với nhiễu}
% Đường accuracy-SNR của ba phương pháp trên cùng hệ trục;
% ngưỡng SNR tối thiểu; ảnh hưởng riêng từng loại nhiễu.
[Nội dung viết tại đây.]

\chohinh{độ chính xác theo SNR của ba phương pháp}

\subsection{Phân tích lỗi}
% Ma trận nhầm lẫn và cách đọc. NHẬN XÉT TRỌNG TÂM: các cặp phím nhầm nhiều nhất
% là cặp cùng hàng hoặc cùng cột. Thống kê khung bị từ chối theo lý do.
% Đánh giá cơ chế chống talk-off.
[Nội dung viết tại đây.]

\subsection{So sánh chi phí tính toán}
% Đối chiếu số phép nhân lý thuyết ở mục 2.2.5 với thời gian chạy đo được;
% giải thích chênh lệch.
[Nội dung viết tại đây.]

\subsection{Nhận xét tổng hợp}
% Bảng so sánh ba phương pháp trên bốn tiêu chí;
% khuyến nghị lựa chọn theo bối cảnh ứng dụng.
[Nội dung viết tại đây.]

\begin{table}[H]
\centering
\caption{So sánh tổng hợp ba phương pháp giải mã}
\label{tab:tong-hop}
\small
\begin{tabular}{l >{\centering\arraybackslash}p{2.3cm} >{\centering\arraybackslash}p{2.3cm} >{\centering\arraybackslash}p{2.3cm} >{\centering\arraybackslash}p{2.4cm}}
\toprule
\thd{Phương pháp} & \thd{Độ chính xác} & \thd{Độ bền nhiễu}
& \thd{Chi phí tính toán} & \thd{Độ phức tạp cài đặt} \\
\midrule
DFT/FFT        & & & & \\
Goertzel       & & & & \\
Ngân hàng lọc  & & & & \\
\bottomrule
\end{tabular}
\end{table}

\clearpage

% ---------------------------------------------------------
% V. KẾT LUẬN VÀ HƯỚNG PHÁT TRIỂN
% ---------------------------------------------------------
\section{KẾT LUẬN VÀ HƯỚNG PHÁT TRIỂN}

\subsection{Kết quả đạt được}
% Đối chiếu trực tiếp với từng mục tiêu ở Mở đầu, kèm số liệu định lượng.
[Nội dung viết tại đây.]

\subsection{Hạn chế}
% Chưa thời gian thực cứng; chưa kiểm chứng trên đường truyền thật;
% dữ liệu chủ yếu tổng hợp; tính năng đã cân nhắc nhưng chưa thực hiện.
[Nội dung viết tại đây.]

\subsection{Hướng phát triển}
% Mở rộng cột 1633 Hz; Goertzel số nguyên trên vi điều khiển;
% chống talk-off nâng cao; tự thích nghi ngưỡng theo nhiễu nền.
[Nội dung viết tại đây.]

\clearpage

% ---------------------------------------------------------
% TÀI LIỆU THAM KHẢO
% ---------------------------------------------------------
\unsection{TÀI LIỆU THAM KHẢO}
\printbibliography[heading=none]

\clearpage

% ---------------------------------------------------------
% PHỤ LỤC
% ---------------------------------------------------------
\begin{appendices}
% Nới ô số trong mục lục để chứa được chữ "PHỤ LỤC n"
\addtocontents{toc}{\protect\addtolength{\protect\cftsecnumwidth}{4.4em}}
\renewcommand{\thesection}{PHỤ LỤC \arabic{section}}
% Hình và bảng trong phụ lục đánh số riêng (PL1.1, PL3.1...) để không
% trùng với hình/bảng của các chương chính.
\renewcommand{\thefigure}{PL\arabic{section}.\arabic{figure}}
\renewcommand{\thetable}{PL\arabic{section}.\arabic{table}}
\setcounter{section}{0}

\section{Mã nguồn các hàm chính}
% dtmf_generate, goertzel_power, dtmf_decode_goertzel,
% design_bpf_bank, dtmf_decide. Dùng lstlisting có đánh số dòng.
\begin{matlabbox}[\texttt{goertzel\_power.m}]
function P = goertzel_power(x, k, N)
%GOERTZEL_POWER Cong suat tai bin k bang thuat toan Goertzel.
    c  = 2*cos(2*pi*k/N);
    s1 = 0; s2 = 0;
    for n = 1:N
        s  = x(n) + c*s1 - s2;
        s2 = s1;
        s1 = s;
    end
    P = s1^2 + s2^2 - c*s1*s2;
end
\end{matlabbox}

\section{Bảng số liệu chi tiết}
% Kết quả đầy đủ theo từng mức SNR, bảng hệ số bộ lọc, bảng chỉ số bin.
[Nội dung viết tại đây.]

\section{Bảng đóng góp của các thành viên}
\begin{table}[H]
\centering
\small
\caption{Đóng góp của từng thành viên}
\label{tab:dong-gop}
\begin{tabular}{@{}c @{\hspace{6pt}}l @{\hspace{6pt}}l p{4.3cm}@{}}
\toprule
\textbf{STT} & \textbf{Họ và tên} & \textbf{Vai trò} & \textbf{Đầu việc cụ thể đã thực hiện} \\
\midrule
1  & Đỗ Hồng Nhất            & Trưởng nhóm, lập trình & \\
2  & Nguyễn Thu Hà           & Điều phối              & \\
3  & Lê Tuấn Anh             & Giao diện              & \\
4  & Nguyễn Quốc Đăng        & Giao diện, đặc tả      & \\
5  & Lê Đức Anh              & Kiểm thử               & \\
6  & Nguyễn Minh Hiếu        & Đặc tả                 & \\
7  & Hồ Văn Hoàng Quân       & Đặc tả                 & \\
8  & Nguyễn Trí Dũng         & Nghiên cứu             & \\
9  & Trương Minh Tân         & Nghiên cứu             & \\
10 & Nguyễn Tiến Mạnh        & Nghiên cứu             & \\
11 & Nguyễn Huy Hoàng        & Nghiên cứu             & \\
12 & Nguyễn Cửu An Phú       & Nghiên cứu             & \\
13 & Đinh Minh Trí           & Nghiên cứu             & \\
14 & Đặng Nguyễn Lâm         & Báo cáo                & \\
15 & Huỳnh Nguyễn Bảo Khương & Báo cáo                & \\
16 & Nguyễn Duy Tiến         & Slide                  & \\
17 & Cao Tiến Thành          & Thuyết trình           & \\
18 & Nguyễn Sỹ Long          & Chưa phân công         & \\
\bottomrule
\end{tabular}
\end{table}

\end{appendices}

\end{document}
